% Chapter: the local seed.  Section 2 of the paper (sec:local-model).
% Owned shared labels: def:local:seed.
% Paper labels: lem:jump-horizontal-mode (Lemma 2.1), prop:local-seed (Proposition 2.2).

Section~2 of the paper builds the seed on \emph{two}-component horizontal fields.  At unit
diffusivity the horizontal equations in the axial Fourier mode $\e^{2\pi ix_3}$ close on
$b\colon\R^2\to\CC^2$, and the weak form of the resulting generator
\eqref{eq:local:horizontal-induction} carries no derivative of the velocity once the
stretching term is integrated by parts.  This is what allows the seed to be built for a
\emph{discontinuous} velocity, the Ponomarenko screw flow, whose unstable eigenvalue is
produced by an intermediate-value argument on an explicit segment of the Bessel dispersion
relation.  The eigenvalue is isolated by the analytic Fredholm theorem, the velocity is
smoothed and the diffusivity is varied in \emph{one} form estimate
\eqref{eq:local:form-small}, and a fixed compactly supported test vector $f_*$ detects the
spectral subspace throughout.  The output is the record of Definition~\ref{def:local:seed},
inhabited at $k_*=1$ by Proposition~\ref{prop:local-seed}; the axial component of the
magnetic field is not part of the seed at all and is recovered only in the finish chapter
(Lemma~\ref{lem:axial-recovery}).

Throughout this chapter $\Lp^2$ inner products are linear in the first argument, as in the
paper; the Lean convention is the opposite one, and every statement below is the
transcription $\inner f\phi=\texttt{inner}\ \CC\ (j\phi)\ f$ of the paper's pairing
(the module docstring of \texttt{FastDynamo/Local/HorizontalWeakForm.lean} records it).

\section{The horizontal operator on $\R^2$ and its weak form}

\begin{definition}[The horizontal generator and its sesquilinear form]
  \label{def:local:horizontal-operator}
  \lean{FastDynamo.HorizontalWeakForm.horizontalDiffusionForm,
    FastDynamo.HorizontalWeakForm.horizontalTransportForm,
    FastDynamo.HorizontalWeakForm.horizontalAxialForm,
    FastDynamo.HorizontalWeakForm.horizontalStretchDivergenceForm,
    FastDynamo.HorizontalWeakForm.horizontalStretchGradientForm,
    FastDynamo.HorizontalWeakForm.horizontalWeakForm,
    FastDynamo.HorizontalWeakForm.HorizontalFormDomain,
    FastDynamo.HorizontalWeakForm.horizontalFormOperator,
    FastDynamo.HorizontalWeakForm.horizontalFormOperator_exists_iff,
    FastDynamo.HorizontalWeakForm.horizontalFormOperator_isEigenpair_iff}
  \leanok
  Fix $N=2\pi$ (the paper writes $4\pi^2$ and $2\pi i$ for $N^2$ and $iN$).  For a velocity $w=(w_\perp,w_3)$ independent of $x_3$ with
  $\Div_\perp w_\perp=0$ the horizontal generator at unit diffusivity is
  \begin{equation}
    A_wb=(\Delta_\perp-N^2)b-w_\perp\cdot\nabla_\perp b-iNw_3b
      +(b\cdot\nabla_\perp)w_\perp,
    \qquad b\colon\R^2\to\CC^2 .
    \label{eq:local:horizontal-induction}
  \end{equation}
  Integrating the stretching term by parts removes every derivative of $w$, so for
  $w\in\Lp^\infty(\R^2;\R^3)$ the sesquilinear form
  \begin{align}
    \mathfrak l_w(b,\phi)
    ={}&\inner{\nabla_\perp b}{\nabla_\perp\phi}+N^2\inner b\phi
       +\inner{w_\perp\cdot\nabla_\perp b}\phi+iN\inner{w_3b}\phi\notag\\
      &+\int_{\R^2}w_\perp\cdot
        \bigl((\Div_\perp b)\overline\phi+(b\cdot\nabla_\perp)\overline\phi\bigr)\dd y
    \label{eq:local:sectorial-form}
  \end{align}
  is defined for $b,\phi\in\Hs^1(\R^2;\CC^2)$.  We write $-A_w$ for the operator on
  $\Lp^2(\R^2;\CC^2)$ associated with $\mathfrak l_w$ in Kato's sense: $b\in D(A_w)$ with
  $A_wb=f$ exactly when $b\in\Hs^1$ and $\mathfrak l_w(b,\phi)=-\inner f\phi$ for every
  $\phi\in\Hs^1$.
\end{definition}

In Lean the five summands of \eqref{eq:local:sectorial-form} are the five declarations
listed above, all on \texttt{WholePlaneC2H1} $=\Hs^1(\R^2;\CC^2)$ with a velocity of type
\texttt{WholePlaneC3LInf}; \texttt{horizontalFormOperator} is the generic
\texttt{associatedFormProdOperator} of \texttt{FastDynamo/Operator/AssociatedForm.lean}
applied to $\mathfrak l_w$ with the pivot $\Hs^1\hookrightarrow\Lp^2$, so the domain is a
graph subspace of $\Hs^1\times\Lp^2$ rather than a subset of $\Lp^2$.  The two
characterisations \texttt{horizontalFormOperator\_exists\_iff} and
\texttt{\_isEigenpair\_iff} are the paper's domain description and the eigenvalue
equation $\mathfrak l_w(b,\phi)=-\lambda\inner b\phi$ verbatim.  The module is
\texttt{FastDynamo/Local/HorizontalWeakForm.lean} and
\texttt{FastDynamo/Local/HorizontalFormOperator.lean}.

\begin{lemma}[Boundedness, Gårding inequality and sectoriality]
  \label{lem:local:garding}
  \lean{FastDynamo.HorizontalWeakForm.norm_horizontalWeakForm_apply_le,
    FastDynamo.HorizontalWeakForm.horizontalWeakForm_garding,
    FastDynamo.HorizontalWeakForm.horizontalWeakForm_abs_im_le,
    FastDynamo.HorizontalWeakForm.horizontalGardingConstant,
    FastDynamo.HorizontalWeakForm.horizontalShiftedWeakForm,
    FastDynamo.HorizontalWeakForm.horizontalShiftedWeakForm_hasComplexCoercivityConstant,
    FastDynamo.HorizontalWeakForm.horizontalShiftedWeakForm_sectorial}
  \leanok
  \uses{def:local:horizontal-operator}
  For $w\in\Lp^\infty(\R^2;\R^3)$ the form $\mathfrak l_w$ is bounded on
  $\Hs^1\times\Hs^1$, and
  \[
    \re\mathfrak l_w(b,b)\geq\tfrac12\norm{\nabla_\perp b}_2^2-C\norm b_2^2,
    \qquad
    \abs{\im\mathfrak l_w(b,b)}\leq C'\bigl(\norm{\nabla_\perp b}_2^2+\norm b_2^2\bigr),
  \]
  with $C,C'$ depending only on $N$ and $\norm w_\infty$.  Consequently, for
  $\alpha=C+\tfrac12$, the shifted form $\mathfrak l_w+\alpha\inner\cdot\cdot$ is coercive on
  $\Hs^1$ with constant $\tfrac12$, and its imaginary part on the diagonal is at most $2C'$
  times its real part; so $\mathfrak l_w$ is closed and sectorial.
\end{lemma}

The Lean constants are explicit: $C=50\norm w_\infty^2+\abs N\norm w_\infty$
(\texttt{horizontalGardingConstant}), $C'=(5+\abs N)\norm w_\infty$, and the sectoriality
slope is $2C'$; all are generous, and no later step uses their value.  The first inequality
is the paper's \texttt{eq:sectorial-estimate-b}.  \textbf{Deviation (shape of the
imaginary-part bound).}  The paper's second display is
$\abs{\im\mathfrak l_w(b,b)}\leq\re\mathfrak l_w(b,b)+C\norm b_2^2$ for every $b\in\Hs^1$,
after increasing $C$ if necessary: a sector of slope one about the shifted real part.  Lean
bounds the imaginary part by the full $\Hs^1$ norm instead and obtains, after the shift by
$\alpha$, a sector of slope $2C'$, which depends on $N$ and $\norm w_\infty$.  Both say that
$\mathfrak l_w$ is sectorial, which is all the paper uses (Kato's representation theorem);
the slope-one form is not formalised.  The coercivity after the shift is stated in the
$\Hs^1$ graph norm rather than in the paper's split form, which is the same statement because
$\norm b_{\Hs^1}^2=\norm b_2^2+\norm{\nabla_\perp b}_2^2$ by construction of the graph
space.

\begin{proof}
  \leanok
  \uses{def:local:horizontal-operator}
  The three velocity terms carrying a derivative are bounded by
  $5\norm w_\infty\norm b_2\bigl(\norm{\partial_1b}_2+\norm{\partial_2b}_2\bigr)$ and the
  axial term by $\abs N\norm w_\infty\norm b_2^2$; Young's inequality
  $5Mm(G_0+G_1)\leq\tfrac14(G_0^2+G_1^2)+50M^2m^2$ absorbs the first into the diffusion
  term and leaves $\tfrac34$ of it, which is more than the claimed $\tfrac12$.  The
  imaginary part is estimated by the same three bounds without absorption, and
  $2\norm{\partial_jb}_2\norm b_2\leq\norm{\partial_jb}_2^2+\norm b_2^2$.  Adding a large
  real multiple of the $\Lp^2$ inner product makes the real part dominate, which is
  coercivity, and the quotient of imaginary by real part is then bounded, which is
  sectoriality.
\end{proof}

\section{An unstable eigenvalue for the discontinuous screw flow}

\begin{definition}[The Ponomarenko screw flow]
  \label{def:local:jump-flow}
  \lean{FastDynamo.HorizontalMode.horizontalJumpInteriorVelocity,
    FastDynamo.HorizontalMode.horizontalJumpVelocity,
    FastDynamo.HorizontalMode.horizontalJumpVelocityLInf,
    FastDynamo.HorizontalMode.horizontalModeVelocity,
    FastDynamo.PonomarenkoFreeOmegaInterface.jumpAxialSpeed,
    FastDynamo.PonomarenkoFreeOmegaInterface.scale_mul_jumpAxialSpeed}
  \leanok
  \uses{def:local:horizontal-operator}
  For $\Omega>0$ put $U=-\Omega/N$ and, in polar coordinates,
  \begin{equation}
    V^0(y)=\one_{\{r<1\}}\bigl(\Omega re_\theta+Ue_3\bigr).
    \label{eq:local:jump-flow}
  \end{equation}
  The normal component vanishes on both sides of $r=1$, so
  $\Div_\perp V^0_\perp=0$ in the sense of distributions and $V^0\in\Lp^\infty$ has
  compact support.
\end{definition}

In Lean $\Omega$ is a \emph{free} real parameter: \texttt{horizontalJumpVelocity Omega U}
is the Cartesian field $(-\Omega y_2,\Omega y_1,U)$ cut off by the indicator of the closed
unit disc, and \texttt{horizontalModeVelocity N Omega} pins $U=-\Omega/N$ through
\texttt{jumpAxialSpeed}.  The paper's open disc $\{r<1\}$ and the Lean closed disc
differ on a null set, hence define the same $\Lp^\infty$ class; this is recorded in the
module docstring.  Making $\Omega$ free is the change that lets the dispersion relation be
read as a \emph{condition on $\kappa$ alone}, which is what the intermediate-value
argument needs.

\begin{definition}[The dispersion relation and its certificate]
  \label{def:local:dispersion}
  \lean{FastDynamo.ponomarenkoDispersion_analyticOnNhd,
    FastDynamo.ModifiedBessel.ponomarenkoDispersion_eq_neg_two_div_mul_deriv_I1_mul_K1,
    FastDynamo.PonomarenkoFreeOmegaInterface.TransverseInterfaceLaws,
    FastDynamo.PonomarenkoFreeOmegaInterface.CertifiedInterfaceData,
    FastDynamo.PonomarenkoFreeOmegaInterface.transverseInterfaceLaws_of_dispersion,
    FastDynamo.ponomarenkoLocalEigenvalue}
  \leanok
  \uses{def:local:jump-flow}
  Separating the angular index $j=1$ and writing $B_\pm=B_r\pm iB_\theta$, the bounded
  solutions of the radial system are the modified Bessel branches $I_2,I_0$ inside and
  $K_2,K_0$ outside, at the transverse wave number $\kappa$ with
  $\lambda_*=\kappa^2-N^2$.  The four transmission conditions
  $\jump{B_r}=\jump{B_\theta}=\jump{B_r'}=0$ and $\jump{B_\theta'}=\Omega B_r(1)$ have a
  non-trivial solution exactly when
  \begin{equation}
    \mathcal D(\kappa):=I_0(\kappa)K_0(\kappa)-I_2(\kappa)K_2(\kappa)
    =-\frac{2i}{\Omega}.
    \label{eq:local:dispersion}
  \end{equation}
  A \emph{certificate} is a triple $(N,\Omega,\kappa)$ with $N>0$, $\re\kappa>0$,
  $\Omega>0$, \eqref{eq:local:dispersion}, the four transmission laws,
  $\re(\kappa^2-N^2)>0$, and the non-vanishing of $I_\nu(\kappa)$ and $K_\nu(\kappa)$ for
  $\nu=0,1,2$ and of the interior radial component at $r=1$.
\end{definition}

\texttt{CertifiedInterfaceData N Omega kappa} is exactly that record; the transmission laws
are a separate record \texttt{TransverseInterfaceLaws}, which
\texttt{transverseInterfaceLaws\_of\_dispersion} derives from
\eqref{eq:local:dispersion} and $\re\kappa>0$, so the record is carried only for
convenience.  The Lean branches are the \emph{cross-normalised} pair
$B_+=K_2(\kappa)I_2(\kappa r)$, $B_-=-K_0(\kappa)I_0(\kappa r)$ inside (and the mirrored
pair outside), a non-zero constant multiple of the paper's solution
(\texttt{eq:bessel-solutions} with $S=C_++C_-=1$) written without denominators; with that
normalisation the two Wronskian jumps
$\jump{B_+'}=-1$, $\jump{B_-'}=+1$ hold identically and
\eqref{eq:local:dispersion} is equivalent to $\Omega B_r(1)=i$.

\begin{lemma}[An admissible root on the segment $\kappa=\varrho(1+it)$]
  \label{lem:local:segment-root}
  \lean{FastDynamo.norm_ponomarenkoSegmentDispersion_sub_inv_cube_le,
    FastDynamo.exists_ponomarenkoSegmentDispersionRe_eq_zero,
    FastDynamo.modifiedBessel_denominators_ne_zero_on_segment,
    FastDynamo.exists_certifiedInterfaceData_of_segment,
    FastDynamo.exists_certifiedInterfaceData_two_pi}
  \leanok
  \uses{def:local:dispersion}
  For $\varrho\geq32$ and $\abs t\leq3/4$,
  \[
    \abs{\varrho^3\mathcal D(\varrho(1+it))-(1+it)^{-3}}\leq\frac{225}{4\varrho^2}
    \leq\frac1{16},
  \]
  and $\re(1+it)^{-3}=(1-3t^2)(1+t^2)^{-3}$ changes sign between $t=\tfrac12$ and
  $t=\tfrac34$ while $\im(1+it)^{-3}\leq-\tfrac13$ there.  Hence there is
  $t_\varrho\in(\tfrac12,\tfrac34)$ with $\re\mathcal D(\varrho(1+it_\varrho))=0$; setting
  $\kappa=\varrho(1+it_\varrho)$ and $\Omega=-2/\im\mathcal D(\kappa)>0$ produces a
  certificate, with
  \[
    \re\lambda_*=\varrho^2(1-t_\varrho^2)-N^2\geq\tfrac7{16}\varrho^2-N^2 .
  \]
  In particular at $N=2\pi$ and $\varrho=32$ a certificate exists with
  $\re\lambda_*\geq448-4\pi^2>0$.
\end{lemma}

\textbf{Deviation (explicit numerals).}  Every threshold here is explicit and is proved in
Lean at the stated numerals; the numerals are choices of the formalisation, not of the
paper.  The paper restricts to $\kappa=\varrho(1+it)$ with $t\in[\tfrac12,\tfrac34]$ and
$\varrho>0$ ``sufficiently large'', uses \texttt{eq:bessel-product-expansion} with an
unquantified $O(\varrho^{-5})$ remainder, and concludes
$\re\lambda_*=\varrho^2(1-t_\varrho^2)-4\pi^2\geq\tfrac7{16}\varrho^2-4\pi^2>0$ after
``increasing $\varrho$ if necessary'' (\texttt{eq:local-eigenvalue}); it fixes no value of
$\varrho$ and no numerical lower bound for $\re\lambda_*$.  Lean proves the expansion with
the explicit error $225/(4\varrho^2)$ for $\varrho\geq32$ ($\varrho=32$ is the smallest power
of two above the binding constraint $\varrho\geq30$), carries the hypothesis
$N^2<\tfrac7{16}\varrho^2$ in place of ``increasing $\varrho$'', and discharges it at
$N=2\pi$, $\varrho=32$ by $(2\pi)^2<448$; hence the numeral $448-4\pi^2$.  The
existence theorem also exposes $\re\kappa=\varrho$, $\im\kappa\in(\varrho/2,3\varrho/4)$
and $\norm\kappa\leq\tfrac54\varrho$, which the later sector estimates consume.

\begin{proof}
  \leanok
  \uses{def:local:dispersion}
  The product $zI_1(z)K_1(z)$ obeys $\abs{zI_1K_1-\tfrac12}\leq3(\re z)^{-2}$ in the
  sector $\norm z\leq2\re z$, by the Gaussian representations of $I_1$ and $K_1$; since
  $(I_1K_1)'(z)=-z\mathcal D(z)/2$, so that $\mathcal D$ is $-2/z$ times the derivative of
  that product, one gets the expansion above on a disc of radius $1/5$
  around each segment point by Cauchy's estimate, and the stated constants follow by
  summing $2\delta+2\eta$.  The sign pattern of $(1+it)^{-3}$ at $t=\tfrac12,\tfrac34$
  survives the error $1/16$, so the intermediate value theorem gives $t_\varrho$.  That the
  four Bessel values $I_0,I_2,K_0,K_2$ do not vanish at $\kappa$ follows from the four
  collapse identities for $z^3\mathcal D(z)$ together with
  $\norm{\kappa^3\mathcal D(\kappa)}\leq2$ and $\norm{\kappa I_1K_1-\tfrac12}<\tfrac12$,
  which force $\norm\kappa\leq3$, contradicting $\re\kappa=\varrho\geq32$.  Finally
  $\im\mathcal D(\kappa)<0$ on the segment makes $\Omega$ positive, and
  $t_\varrho^2\leq9/16$ gives the growth floor.
\end{proof}

\begin{definition}[The Cartesian eigenmode]
  \label{def:local:cartesian-mode}
  \lean{FastDynamo.HorizontalMode.horizontalMode,
    FastDynamo.HorizontalMode.horizontalEigenvalue,
    FastDynamo.HorizontalMode.horizontalMode_ne_zero,
    FastDynamo.HorizontalMode.horizontalMode_tail_le,
    FastDynamo.HorizontalMode.HorizontalModeWeakEquation,
    FastDynamo.HorizontalMode.horizontalModeWeakEquation_iff_distributional}
  \leanok
  \uses{def:local:dispersion, def:local:horizontal-operator}
  Given a certificate, the field $b$ with
  \[
    b_1+ib_2=\e^{2i\theta}B_+(r),\qquad b_1-ib_2=B_-(r)
  \]
  is smooth on each side of $r=1$, continuous across it, decays exponentially, and lies in
  $\Hs^1(\R^2;\CC^2)$.  Its eigenvalue is $\lambda_*=\kappa^2-N^2$.
\end{definition}

The Lean mode is built branch by branch: the origin regularity comes from the
$r^2$-analyticity of $I_0$ and $I_2$, the exterior decay from the Gaussian representation
of $K_\nu$, and the $\Hs^1$ membership from a slice fundamental theorem of calculus.  The
weak eigenvalue equation is stated in two equivalent forms, against all
$\phi\in\Hs^1$ and as an identity of tempered distributions; the equivalence rests on the
injectivity of $\Hs^{-1}(\R^2;\CC^2)\to\cS'(\R^2;\CC^2)$ and is a theorem, not a
definitional move.

\begin{lemma}[An unstable eigenvalue of the discontinuous flow]
  \label{lem:jump-horizontal-mode}
  \lean{FastDynamo.HorizontalMode.horizontalModeWeakIdentity_holds,
    FastDynamo.HorizontalMode.horizontalMode_isEigenpair_of_certified,
    FastDynamo.HorizontalMode.exists_horizontalMode_eigenpair_two_pi}
  \leanok
  \uses{def:local:jump-flow, def:local:cartesian-mode, lem:local:segment-root,
    def:local:horizontal-operator}
  There exists $\Omega>0$ such that, for the velocity $V^0$ of
  \eqref{eq:local:jump-flow}, the operator $A_{V^0}$ has an eigenvalue $\lambda_*$ with
  $\re\lambda_*>0$.  An associated non-zero eigenfunction lies in the angular mode
  $b=\e^{i\theta}(B_re_r+B_\theta e_\theta)$, is smooth on each side of $r=1$ and decays
  exponentially at infinity.
\end{lemma}

\textbf{Deviation: cross-normalised branches instead of the pair $(C_+,C_-)$.}  The
paper writes the interior and exterior solutions as
$C_\pm I_{\nu_\pm}(\kappa r)/I_{\nu_\pm}(\kappa)$ and
$C_\pm K_{\nu_\pm}(\kappa r)/K_{\nu_\pm}(\kappa)$ (\texttt{eq:bessel-solutions}), and
therefore needs $I_\nu(\kappa)K_\nu(\kappa)\neq0$ for its denominators; it obtains this
from the expansion \texttt{eq:bessel-product-expansion},
$I_\nu(\kappa)K_\nu(\kappa)=\tfrac1{2\kappa}-\tfrac{4\nu^2-1}{16\kappa^3}+O(\varrho^{-5})$
for $\nu\in\{0,2\}$, for $\varrho$ sufficiently large, and then normalises by
$S=C_++C_-=1$.  Lean cross-normalises instead -- $K_\nu(\kappa)I_\nu(\kappa r)$ inside and
$I_\nu(\kappa)K_\nu(\kappa r)$ outside -- so the two branches match at $r=1$ identically,
the Wronskian gives the derivative jumps $\mp1$ with no denominator, and the eigenfunction
obtained is a non-zero constant multiple of the paper's eigenfunction with $S=1$.  The
Bessel values do not vanish: $I_\nu(\kappa),K_\nu(\kappa)\neq0$
for $\nu=0,1,2$ on the whole segment for $\varrho\geq32$
(\texttt{modifiedBessel\_denominators\_ne\_zero\_on\_segment}, by the argument given in the
proof of Lemma~\ref{lem:local:segment-root} rather than from the expansion); these facts are
fields of the certificate, and $K_0(\kappa)\neq0$ is how $b\neq0$ is read off, the first
Cartesian component of the Lean field at the origin being $-K_0(\kappa)/2$.  Lean also
routes the expansion through the single order-one product $zI_1(z)K_1(z)$, using the exact
algebraic identity $z^3\mathcal D(z)=2z-4zI_1(z)K_1(z)$, rather than expanding $\nu=0$ and
$\nu=2$ separately: the sectorised estimate $\abs{zI_1(z)K_1(z)-\tfrac12}\leq3/(\re z)^2$
(\texttt{FastDynamo/\allowbreak Local/\allowbreak PonomarenkoSectorOrderOneProduct.lean}) is the effective form of
\texttt{eq:bessel-product-expansion}, and the segment statement
$\varrho^3\mathcal D(\varrho(1+it))=(1+it)^{-3}+O(\varrho^{-2})$ carries the explicit
constant $225/(4\varrho^2)$.  The root itself is produced, in the paper and here, by the
intermediate value theorem in $t\in[\tfrac12,\tfrac34]$.

At $N=2\pi$ the Lean theorem has no hypothesis and also returns the numeral
$448-4\pi^2\leq\re\lambda_*$, which comes from the formalisation's choice $\varrho=32$
(Lemma~\ref{lem:local:segment-root}); the paper asserts only $\re\lambda_*>0$, via
$\re\lambda_*\geq\tfrac7{16}\varrho^2-4\pi^2$.  The auxiliary clauses of the paper's
statement -- the angular form of the eigenfunction, smoothness on each side of the circle,
exponential decay -- are separate declarations (\texttt{horizontalMode},
\texttt{horizontalInteriorField\_contDiff}, \texttt{horizontalMode\_tail\_le}) rather than
conjuncts of the headline theorem.  The paper's sentence ``the transmission conditions
cancel the interface distributions'' is, in Lean, a \emph{pointwise} identity between the
two one-sided jump densities on the unit circle, so the result does not depend on the
normalisation of the circle measure.

\begin{proof}
  \leanok
  \uses{def:local:cartesian-mode, lem:local:segment-root, def:local:jump-flow}
  Writing $D_j=\partial_jb+w_\perp b_j$ for the pivot field, the weak equation
  $\mathfrak l_{V^0}(b,\phi)=-\lambda_*\inner b\phi$ is equivalent to the distributional
  identity $\sum_j\partial_jD_j=M+\lambda_*b$ with
  $M=N^2b+w_\perp\cdot\nabla_\perp b+iNw_3b+w_\perp\Div_\perp b$.  Away from $r=1$ this is
  the bulk identity: the Cartesian Laplacian acts on the two circular combinations
  $b_1\pm ib_2$ with angular indices $(j\pm1)^2=4,0$, giving $\Delta_\perp b=\kappa^2b$,
  while the
  transport and stretching terms cancel because $(b\cdot\nabla_\perp)V^0_\perp
  =V^0_\perp\cdot\nabla_\perp b-i\Omega b$ and $iNU=-i\Omega$.  Across $r=1$ the
  distributional derivative of $D_j$ acquires a surface term with density
  $n_j\jump{D_j}$; the outward normal is the position vector and
  $\jump{\partial_rb}=\Omega B_re_\theta$ is the fourth transmission law, which is
  \eqref{eq:local:dispersion}.  The two one-sided densities are therefore pointwise
  opposite and the surface term vanishes.  Lemma~\ref{lem:local:segment-root} supplies a
  certificate, and $b\neq0$ because its first component at the origin is
  $-K_0(\kappa)/2\neq0$.
\end{proof}

\section{Spectral isolation and regularisation}

\begin{lemma}[Compactness of the velocity form and the analytic Fredholm family]
  \label{lem:local:compactness}
  \lean{FastDynamo.HorizontalWeakForm.freeHorizontalForm,
    FastDynamo.HorizontalWeakForm.horizontalVelocityForm,
    FastDynamo.HorizontalWeakForm.horizontalPencilForm,
    FastDynamo.HorizontalWeakForm.freeHorizontalForm_hasComplexCoercivityConstant,
    FastDynamo.HorizontalWeakForm.freeHorizontalForm_bijective,
    FastDynamo.HorizontalWeakForm.isCompactOperator_horizontalVelocityZeroOrder_of_ae_zero_off,
    FastDynamo.isCompactOperator_adjoint}
  \leanok
  \uses{def:local:horizontal-operator, lem:local:garding}
  Let $w\in\Lp^\infty(\R^2;\R^3)$ vanish almost everywhere off a set of finite measure.
  On $V=\Hs^1(\R^2;\CC^2)$ with antidual $V^*=\Hs^{-1}$, write
  $\mathbb A_w(z)=F_z+K_w$ with $F_z=z-(\Delta_\perp-N^2)$.  Then $F_z\colon V\to V^*$ is
  an isomorphism for $\re z>0$, $K_w\colon V\to V^*$ is compact, and
  $\mathbb A_w(z)=F_z(\id+F_z^{-1}K_w)$ is an analytic Fredholm family on $\{\re z>0\}$.
\end{lemma}

Three deviations.  First, the hypothesis on the velocity: the paper takes $w$ with compact
support, $\supp w\subset D(0,R_w)$, and restricts to that disc; Lean assumes only that $w$
vanishes almost everywhere off a measurable set $s$ of finite measure and restricts to $s$,
local Rellich compactness being proved for finite-measure sets
(\texttt{FastDynamo/Analysis/WholePlaneBandLimitedRellichUnconditional.lean}).  This is more
general, and is applied to $V^0$, which vanishes off the unit disc.  Second, the paper
inverts $F_z$ by the Fourier multiplier $(z+\abs\xi^2+4\pi^2)^{-1}$; Lean instead observes that
$\re F_z(b,b)=\norm{\nabla b}_2^2+(N^2+\re z)\norm b_2^2$ and applies complex
Lax--Milgram, with coercivity constant $\min(1,\re z)$.  A form is determined by its
values, so the operator is the same.  Third, Mathlib has no Schauder theorem, so the
compactness of the adjoint $R^*$ needed for the ``derivative on the input'' half of the
argument is \texttt{isCompactOperator\_adjoint}, proved here from scratch in the
Hilbert-space shortcut $\norm{S^\dagger y}^2=\re\inner y{SS^\dagger y}$.  The argument itself
is the paper's: §2.3 of the paper takes $R\colon\Hs^1(\R^2)\to\Lp^2(D(0,R_w))$ compact by
Rellich, $R^*\colon\Lp^2(D(0,R_w))\to\Hs^{-1}(\R^2)$ compact as its adjoint, and decomposes
each of the three kinds of term through $R$, multiplication and $R^*$; only the proof of the
adjoint step is supplied here.

\begin{proof}
  \leanok
  \uses{lem:local:garding}
  Let $s$ be the finite-measure set off which $w$ vanishes.  Terms with a derivative on
  the test function factor through the compact restriction $\Hs^1(\R^2)\to\Lp^2(s)$;
  terms with a derivative on the input factor through its adjoint, a compact map
  $\Lp^2(s)\to\Hs^{-1}(\R^2)$, after the coefficient has been moved across the
  inner product (multiplication by a real $\Lp^\infty$ coefficient is self-adjoint); the
  zeroth-order term factors through the same restriction.  Multiplication by a bounded
  coefficient is continuous in each case, so $K_w$ is compact by Rellich's theorem.  Since
  $F_z$ is affine in $z$, $z\mapsto F_1^{-1}F_z$ is entire and its inverse is analytic on
  the invertible locus by the Neumann series, which makes the family analytic Fredholm.
\end{proof}

\begin{lemma}[Discreteness of the right-half-plane spectrum and the isolating contour]
  \label{lem:local:isolation}
  \lean{FastDynamo.HorizontalWeakForm.horizontalSpectrum,
    FastDynamo.HorizontalWeakForm.isDiscrete_horizontalSpectrum,
    FastDynamo.HorizontalWeakForm.exists_ne_zero_horizontalPencilForm_eq_zero,
    FastDynamo.HorizontalWeakForm.exists_isEigenpair_of_mem_horizontalSpectrum,
    FastDynamo.HorizontalWeakForm.mem_horizontalSpectrum_iff,
    FastDynamo.HorizontalWeakForm.exists_punctured_ball_isInvertible_horizontalPencilForm,
    FastDynamo.HorizontalWeakForm.HorizontalIsolation,
    FastDynamo.HorizontalWeakForm.HorizontalIsolation.projection,
    FastDynamo.HorizontalWeakForm.HorizontalIsolation.isIdempotentElem_projection,
    FastDynamo.HorizontalWeakForm.HorizontalIsolation.finiteDimensional_projection_range,
    FastDynamo.HorizontalWeakForm.HorizontalIsolation.projection_apply_eigenvector,
    FastDynamo.HorizontalWeakForm.HorizontalIsolation.projection_ne_zero,
    FastDynamo.HorizontalMode.modeIsolation,
    FastDynamo.HorizontalMode.modeIsolation_spec}
  \leanok
  \uses{lem:local:compactness, lem:jump-horizontal-mode}
  Let $w$ be as in Lemma~\ref{lem:local:compactness}.  The set of $z$ with $\re z>0$ at
  which $\mathbb A_w(z)$ is not invertible is discrete, every point of $\{\re z>0\}$ has a
  punctured neighbourhood on which $\mathbb A_w(z)$ is invertible, and that set is exactly the
  set of eigenvalues of $A_w$ with positive real part; so the spectrum of $A_w$ in the open
  right half-plane consists of isolated eigenvalues.  Applied to $w=V^0$ this fixes a
  positively oriented circle $\Gamma$, centred at $\lambda_*$ and of radius at most
  $\tfrac14\re\lambda_*$, whose closed interior $\cG$ lies in $\{\re z>0\}$ and meets that
  set only in $\lambda_*$; the Riesz projection
  $P^0=(2\pi i)^{-1}\int_\Gamma(z-A_{V^0})^{-1}\dd z$ is idempotent, of finite rank, and
  satisfies $P^0b=b\neq0$ for the eigenmode $b$.
\end{lemma}

The paper's analytic Fredholm paragraph in §2.3 concludes that $\mathbb A_w(z)$ is invertible
except at a discrete set of points of the open right half-plane, that at each exceptional point
Fredholm index zero gives a non-zero kernel, which the definition of $A_w$ identifies with
$\ker(z-A_w)$, and hence that the spectrum of $A_w$ in the open right half-plane consists of
isolated eigenvalues.  Lean formalises all of it: the discreteness
(\texttt{isDiscrete\_horizontalSpectrum}, where \texttt{horizontalSpectrum} is \emph{defined}
as the non-invertibility set of $\mathbb A_w(z)$ in $\{\re z>0\}$), the punctured
neighbourhoods, and the identification of that set with the eigenvalues
(\texttt{mem\_horizontalSpectrum\_iff}, \texttt{FastDynamo/Local/HorizontalSpectrumEigenvalue.lean}).
The non-zero kernel at an exceptional point comes from the factorisation
$\mathbb A_w(z)=F_z(\id-K(z))$ and the Fredholm alternative for the compact $K(z)$, which is
the paper's index-zero statement in the form ``injective implies invertible''.  The eigenvector
lies in the form-operator domain of $A_w$, which is the paper's $D(A_w)$ for $w\in\Lp^\infty$ with
compact support; the hypotheses are those of \texttt{isDiscrete\_horizontalSpectrum}
($w$ vanishing off a set of finite measure).  In the other direction, Lean proves that the
Riesz projection of every isolating circle has finite rank
(\texttt{HorizontalIsolation.finiteDimensional\_projection\_range}), which the paper does not
state.  The radius bound $r\leq\tfrac14\re\lambda_*$, which gives
$g_*\geq\tfrac34\re\lambda_*$, is a choice of the formalisation; the paper chooses $\Gamma$
only qualitatively.

The Lean contour is a \texttt{SpectralCircle}: a centre, a radius, and the proof that the
closed disc lies in $\{\re z>0\}$; its \texttt{growthFloor} is
$\re(\text{centre})-\text{radius}$, which \emph{is} $\min_{\zeta\in\cG}\re\zeta$.  The
resolvent delivered by the isolation bundle is domain valued, that is $\Lp^2\to D(A_w)$
with the $\Hs^1$ graph domain; the upgrade to $\Lp^2\to\Hs^2$ is
Lemma~\ref{lem:local:h2-resolvent} and is where \eqref{eq:local:continuation-bound} is
proved.

\begin{proof}
  \leanok
  \uses{lem:local:compactness, lem:jump-horizontal-mode}
  The analytic Fredholm theorem applied to $\id+F_z^{-1}K_w$ on the connected open set
  $\{\re z>0\}$, at the invertible real base point supplied by the Gårding estimate of
  Lemma~\ref{lem:local:garding} (the paper uses \texttt{eq:sectorial-estimate-b} in the
  same way), shows that the non-invertibility set is discrete.  An eigenvector of $A_w$
  makes $\mathbb A_w$ singular at its eigenvalue; conversely, at a point of non-invertibility
  the Fredholm alternative for the compact part gives a non-zero null vector, which solves the
  weak eigenvalue identity and is therefore an eigenvector.  Around $\lambda_*$ pick a punctured ball
  of invertibility; a circle of radius at most $\tfrac14\re\lambda_*$ inside it is $\Gamma$.
  Idempotence of the contour integral is the annular resolvent identity, finite rank
  follows from the compactness of the state resolvent, and $P^0b=b$ is Cauchy's formula
  for $(z-A_{V^0})^{-1}b=(z-\lambda_*)^{-1}b$.
\end{proof}

\begin{definition}[The smoothed velocity and the diffusivity family]
  \label{def:local:smoothed-velocity}
  \lean{FastDynamo.HorizontalMode.horizontalSmoothedVelocityLInf,
    FastDynamo.HorizontalMode.horizontalSmoothedModeVelocity,
    FastDynamo.HorizontalMode.horizontalSmoothingErrorL4,
    FastDynamo.HorizontalWeakForm.horizontalWeakFormDelta,
    FastDynamo.HorizontalWeakForm.horizontalPencilFormDelta,
    FastDynamo.HorizontalWeakForm.horizontalFormDeltaOperator,
    FastDynamo.horizontalLocalOperator}
  \leanok
  \uses{def:local:jump-flow, def:local:horizontal-operator}
  For $0<h<1/2$ let $\chi_h$ be a smooth radial cut-off with $0\leq\chi_h\leq1$, equal to
  one on $[0,1-h]$ and zero on $[1+h,\infty)$, and set
  \begin{equation}
    v_h(y)=\chi_h(\abs y)(-\Omega y_2,\Omega y_1,U)\in\C_c^\infty(\R^2;\R^3),
    \qquad\Div_\perp(v_h)_\perp=0 .
    \label{eq:local:smooth-velocity}
  \end{equation}
  This display is the paper's \texttt{eq:smooth-local-velocity}; the
  solenoidality is immediate from $\nabla\chi_h(\abs y)\cdot y^\perp=0$.
  For $\delta>0$ put
  $\mathfrak l_{h,\delta}=\mathfrak l_{v_h}
   +(\delta-1)\bigl(\inner{\nabla_\perp\cdot}{\nabla_\perp\cdot}+N^2\inner\cdot\cdot\bigr)$,
  with shifted form operator $\mathbb A_{h,\delta}(z)$ and generator
  \begin{equation}
    \cL_\delta b=\delta(\Delta_\perp-N^2)b-v_\perp\cdot\nabla_\perp b-iNv_3b
      +(b\cdot\nabla_\perp)v_\perp,
    \qquad D(\cL_\delta)=\Hs^2(\R^2;\CC^2).
    \label{eq:local:generator}
  \end{equation}
\end{definition}

A note on $N$.  The paper has no axial wave number parameter: it works in the axial mode
$\e^{2\pi ix_3}$ and writes $4\pi^2$ and $2\pi i$ in \texttt{eq:horizontal-induction},
\texttt{eq:local-generator} and everywhere else.  This chapter keeps the symbol $N$, because
the Lean
development does: \texttt{LocalSeedParameters.wavenumber} is $2\pi k_*$ with $k_*$ a
positive natural-number field, and the concrete inhabitant is taken at $k_*=1$
(\texttt{exists\_localSeedC2\_two\_pi}), so every statement below is the paper's at
$k_*=1$ and is more general otherwise.  Keeping $k_*$ free costs nothing and is what makes
the rescaling (Lemma~5.2 of the paper) in a later chapter a statement about an arbitrary scale family.

The whole $\delta$-dependence is the single affine term $(\delta-1)\cdot$(diffusion form),
which is what makes the constant of \eqref{eq:local:form-small} independent of $z$: a
difference of two members of the family is $z$-free, the free form $F_z$ cancelling.
\texttt{horizontalLocalOperator N delta w} is the \emph{strong} operator
\eqref{eq:local:generator} on $\Hs^2$; that it agrees with the form operator on the form
domain is \texttt{FastDynamo/Local/HorizontalFormStrongIdentification.lean} and is the
bridge the resolvent bound below needs.

\begin{lemma}[The one form perturbation estimate]
  \label{lem:local:form-small}
  \lean{FastDynamo.HorizontalWeakForm.horizontalPencilFormDelta_sub,
    FastDynamo.HorizontalWeakForm.complexFormOpNorm_horizontalPencilFormDelta_sub_le,
    FastDynamo.HorizontalMode.exists_norm_horizontalSmoothingErrorL4_le,
    FastDynamo.HorizontalWeakForm.exists_complexFormOpNorm_horizontalModePencil_sub_le}
  \leanok
  \uses{def:local:smoothed-velocity, def:local:jump-flow}
  There is $C$, independent of $z$, of $h\in(0,\tfrac12)$ and of $\delta$, with
  \begin{equation}
    \norm{\mathbb A_{h,\delta}(z)-\mathbb A_{0,1}(z)}_{V\to V^*}
    \leq C\bigl(h^{1/4}+\abs{\delta-1}\bigr).
    \label{eq:local:form-small}
  \end{equation}
\end{lemma}

\begin{proof}
  \leanok
  \uses{def:local:smoothed-velocity}
  The difference splits as $(K_{v_h}-K_{V^0})+(\delta-1)\cdot$(diffusion form) with no $z$
  in it.  The velocity difference is uniformly bounded and supported in an annulus of area
  $O(h)$, so $\norm{v_h-V^0}_4\leq Ch^{1/4}$; each first-order term is then estimated with
  H\"older exponents $(4,4,2)$ or $(4,2,4)$ according to which field carries the
  derivative, using $\Hs^1(\R^2)\hookrightarrow\Lp^4(\R^2)$, and the coefficient is moved
  onto the test function by self-adjointness of multiplication by a real $\Lp^\infty$
  field.  The zeroth-order term is bounded the same way and the diffusion difference
  directly on $V\times V$.
\end{proof}

\begin{lemma}[The Riesz projections move by less than a quarter]
  \label{lem:local:riesz-perturbation}
  \lean{FastDynamo.HorizontalWeakForm.HorizontalNeumannData,
    FastDynamo.HorizontalWeakForm.HorizontalNeumannData.projection,
    FastDynamo.HorizontalWeakForm.HorizontalNeumannData.norm_stateResolvent_sub_le_two,
    FastDynamo.HorizontalWeakForm.HorizontalNeumannData.norm_projection_sub_le_two,
    FastDynamo.HorizontalMode.exists_seedNeumannData,
    FastDynamo.HorizontalMode.exists_seedThresholds}
  \leanok
  \uses{lem:local:form-small, lem:local:isolation}
  The inverses $\mathbb A_{0,1}(z)^{-1}$ are uniformly bounded on $\Gamma$, so for
  $h^{1/4}+\abs{\delta-1}$ small a Neumann series gives
  $\sup_{z\in\Gamma}\norm{\mathbb A_{h,\delta}(z)^{-1}-\mathbb A_{0,1}(z)^{-1}}
  \leq C(h^{1/4}+\abs{\delta-1})$; restricting to $\Lp^2$ gives the resolvents of
  $\cL_{h,\delta}$ on $\Gamma$, and integrating gives
  $\norm{P_{h,\delta}-P^0}_{2\to2}\leq C(h^{1/4}+\abs{\delta-1})$.  There are thresholds
  $h_1,\eta_1>0$ with $\norm{P_{h,\delta}-P^0}<\tfrac14$ whenever $h\leq h_1$ and
  $\abs{\delta-1}\leq\eta_1$.
\end{lemma}

The Lean perturbation is run on the open annulus $B(\lambda_*,3r/2)\setminus\bar
B(\lambda_*,r/4)$ rather than on $\Gamma$ alone; this costs nothing, because
\eqref{eq:local:form-small} is $z$-free, and it buys idempotence of the perturbed
projection from the annular Riesz calculus.  The explicit thresholds are
$h_1=\min(h_0,(8(K+1))^{-4})$ and $\eta_1=\min(\eta_0,(8(K+1))^{-1})$, which give
$K(h^{1/4}+\abs{\delta-1})\leq K/(4(K+1))<\tfrac14$.

\begin{proof}
  \leanok
  \uses{lem:local:form-small, lem:local:isolation}
  Uniform boundedness of $\mathbb A_{0,1}(z)^{-1}$ on the compact annulus is continuity of
  the inverse plus compactness.  With $\eta=C(h^{1/4}+\abs{\delta-1})$ bounding the
  perturbation, the algebraic identity $S^{-1}=\id-(TD)S^{-1}$ gives
  $\norm{S^{-1}}\leq(1-B\eta)^{-1}$ and hence the resolvent difference; the contour
  integral of a bound is the bound times the length of $\Gamma$ divided by $2\pi$, that is
  times the radius.  The thresholds are the stated explicit minima.
\end{proof}

\begin{lemma}[A compactly supported spectral witness]
  \label{lem:local:witness}
  \lean{FastDynamo.rayleighRe,
    FastDynamo.wholePlaneC2DiskCutoff,
    FastDynamo.exists_diskSupported_witness,
    FastDynamo.HorizontalMode.exists_horizontalSeedWitness}
  \leanok
  \uses{lem:local:isolation}
  There are $L_*>0$ and $f_*\in\Lp^2(\R^2;\CC^2)$ with $\norm{f_*}_2=1$, $f_*=0$ almost
  everywhere off $D(0,L_*)$, and $\re\inner{P^0f_*}{f_*}>\tfrac34$.
\end{lemma}

\textbf{Deviation.}  The paper asks for $f_*\in\C_c^\infty(\R^2;\CC^2)$ in Proposition~2.2
and obtains it by density in the proof; Lean produces
the indicator truncation $\one_{D(0,L)}x$ of a normalised element of $\Ran P^0$,
renormalised, which is an $\Lp^2$ field with compact support and not a smooth one.  Nothing
downstream uses smoothness: the global layer consumes only $\norm{f_*}_2=1$, the almost
everywhere support statement and $\re\inner{P_\delta f_*}{f_*}>\tfrac12$.  The support
statement is also realised on the \emph{closed} disc ($L_*<\norm y$ implies $f_*(y)=0$
a.e.), which is the weaker and therefore the safe direction, since
\eqref{eq:transfer:scaled-witness-support} is strict.

\begin{proof}
  \leanok
  \uses{lem:local:isolation}
  Take a unit $x\in\Ran P^0$, which exists because $P^0\neq0$.  Then
  $\re\inner{P^0x}x=1$, and $\one_{D(0,L)}x\to x$ in $\Lp^2$ as $L\to\infty$ by dominated
  convergence, so the normalised Rayleigh quotient
  $y\mapsto\re\inner{P^0y}y/\norm y^2$, which is continuous away from the origin, exceeds
  $\tfrac34$ for $L$ large.
\end{proof}

\begin{lemma}[The uniform $\Lp^2\to\Hs^2$ resolvent bound]
  \label{lem:local:h2-resolvent}
  \lean{FastDynamo.exists_wholePlaneH2_of_laplacian_eq_l2,
    FastDynamo.HorizontalWeakForm.horizontalStrongPairMap,
    FastDynamo.HorizontalWeakForm.norm_horizontalStrongPairMap_le_bound,
    FastDynamo.HorizontalWeakForm.horizontalLocalResolventBoundOfBounds,
    FastDynamo.HorizontalWeakForm.horizontalStrongResolvent_stateResolvent,
    FastDynamo.HorizontalWeakForm.horizontalStrongResolvent_rieszProjection}
  \leanok
  \uses{def:local:smoothed-velocity, lem:local:riesz-perturbation}
  For $\delta$ in a compact interval $I=[\delta_-,\delta_+]$ with $\delta_->0$ the form
  domain of $\mathfrak l_{h,\delta}$ is $\Hs^2(\R^2;\CC^2)$, the strong operator
  $\cL_\delta$ of \eqref{eq:local:generator} has the same resolvent as the form operator,
  hence the same Riesz projection, and there is one constant $L$ with
  \begin{equation}
    \sup_{\delta\in I,\ z\in\Gamma}
      \norm{(z-\cL_\delta)^{-1}}_{\Lp^2\to\Hs^2}\leq L .
    \label{eq:local:continuation-bound}
  \end{equation}
\end{lemma}

\textbf{Lean proves more than the paper asks.}  The paper's
\texttt{eq:local-continuation-bound} in Proposition~2.2 (\texttt{prop:local-seed}) is the
$\Lp^2\to\Hs^1$ bound; what is stated and proved here is the $\Lp^2\to\Hs^2$ bound, and the
paper's bound is its corollary through the embedding $\Hs^2\hookrightarrow\Hs^1$.  This is not extra work.  The Lean
$\cL_\delta$ is a \texttt{DomainOperator} with bundled domain $\Hs^2(\R^2;\CC^2)$ -- the
domain identification that the paper's proof of Proposition~2.2 performs too -- and
\texttt{HasResolventOn} stores a domain-valued inverse $\Lp^2\to\Hs^2$, so the $\Hs^2$
operator norm is the only norm available at that type.  Only the $\Lp^2$ and gradient halves
are ever consumed, by \texttt{norm\_embedding\_inv\_apply\_le} and
\texttt{norm\_deriv\_inv\_apply\_le}, which is exactly the $\Hs^1$ content that
Corollary~\ref{cor:scaled-local-resolvents} needs.

The $\Hs^2$ estimate $\norm b_{\Hs^2}\leq C(\norm b_2+\norm{\Delta_\perp b}_2)$ is obtained
from Banach's isomorphism theorem applied to $h\mapsto(h,\Delta h)$ onto the closed pair
space, not from the Fourier symbol; the constant is therefore unspecified, which is all
\eqref{eq:local:continuation-bound} asks for.  The identification of the two Riesz
projections is the step that transports the witness clause, proved for the form operator,
to the strong generator that the global layer uses.

\begin{proof}
  \leanok
  \uses{def:local:smoothed-velocity, lem:local:riesz-perturbation}
  An element of the form-operator domain lies in $\Hs^1$; since $v_h$ is smooth with
  bounded derivatives the lower-order terms lie in $\Lp^2$, so
  $\delta\Delta_\perp b\in\Lp^2$ and $b\in\Hs^2$, with
  $\delta\norm{\Delta_\perp b}_2\leq\norm f_2+C_v\norm b_{\Hs^1}$.  The form inverse gives
  $\norm b_{\Hs^1}\leq C\norm f_2$ uniformly on $I\times\Gamma$ by
  Lemma~\ref{lem:local:riesz-perturbation}, and $\delta\geq\delta_->0$ turns the two into
  one $\delta$-free bound.  The two resolvents have the same ambient value because the
  $\Hs^2$ lift of a graph element has the same $\Lp^2$ value, and Riesz projections of
  operators with equal state resolvents on the same circle agree.
\end{proof}

\section{The local seed package}

\begin{definition}[The local seed]
  \label{def:local:seed}
  \lean{FastDynamo.LocalSeedParameters, FastDynamo.LocalSeedC2,
    FastDynamo.LocalSeedC2.velocity,
    FastDynamo.LocalSeedC2.coefficients,
    FastDynamo.LocalSeedC2.resolver,
    FastDynamo.LocalSeedC2.norm_resolver_le,
    FastDynamo.LocalSeedC2.projection,
    FastDynamo.LocalSeedC2.re_inner_projection_witness_gt,
    FastDynamo.LocalSeedC2.growthFloor,
    FastDynamo.LocalSeedC2.growthFloor_pos}
  \leanok
  \uses{def:local:smoothed-velocity, lem:local:h2-resolvent, lem:local:witness}
  A \emph{local seed} is the data
  \[
    \bigl(N=2\pi k_*,\ v=v_h,\ I=[\delta_-,\delta_+],\ \Gamma,\ \cG,\ g_*,\ f_*,\ L_*,\ L\bigr)
  \]
  with $0<\delta_-<1<\delta_+$, $v$ helical with $\Div_\perp v_\perp=0$, $\Gamma$ a
  positively oriented circle with $\cG\subset\{\re z>0\}$,
  $\Gamma\subset\rho(\cL_\delta)$ for every $\delta\in I$ with the uniform bound
  \eqref{eq:local:continuation-bound}, the witness clause
  \eqref{eq:local:spectral-witness} and the gap
  $g_*=\min_{\zeta\in\cG}\re\zeta>0$.
\end{definition}

The Lean record is \texttt{FastDynamo.LocalSeedC2}
(\texttt{FastDynamo/Local/LocalSeedC2.lean}), extending
\texttt{LocalSeedParameters} by the velocity data, the family of resolvents and the
witness.  The velocity is a \emph{definition} of the record, not a field, so the record
cannot be inhabited by a non-helical field, and $\Div_\perp v_\perp=0$ is a theorem about
it.  The growth floor is \texttt{contour.growthFloor}$=\re(\text{centre})-\text{radius}$,
which is exactly $\min_{\zeta\in\cG}\re\zeta$ over the closed disc.

\begin{proposition}[Uniform instability of the reference model on $\R^2$]
  \label{prop:local-seed}
  \lean{FastDynamo.HorizontalMode.exists_localSeedC2_two_pi,
    FastDynamo.HorizontalMode.localSeedC2,
    FastDynamo.HorizontalMode.localSeedC2_kStar,
    FastDynamo.HorizontalMode.localSeedC2_wavenumber,
    FastDynamo.HorizontalMode.seedGrowthFloorBound,
    FastDynamo.HorizontalMode.seedGrowthFloorBound_le_localSeedC2_growthFloor,
    FastDynamo.HorizontalMode.localSeedC2_angularSpeed_pos}
  \leanok
  \uses{def:local:seed, lem:local:riesz-perturbation, lem:local:witness,
    lem:local:h2-resolvent, lem:jump-horizontal-mode}
  With $N=2\pi$ there are $h\in(0,\tfrac12]$, an interval $I=[\delta_-,\delta_+]$ with
  $0<\delta_-<1<\delta_+$, a compactly supported unit $f_*$ and $L_*>0$ with the following
  properties.  For the velocity $v=v_h$ of \eqref{eq:local:smooth-velocity} and
  $\delta\in I$, the generator $\cL_\delta$ of \eqref{eq:local:generator} satisfies
  $\Gamma\subset\rho(\cL_\delta)$ together with \eqref{eq:local:continuation-bound}, and the
  Riesz projections $P_\delta=(2\pi i)^{-1}\int_\Gamma(z-\cL_\delta)^{-1}\dd z$ obey
  \begin{equation}
    \norm{f_*}_2=1,\qquad\supp f_*\subset D(0,L_*),\qquad
    \re\inner{P_\delta f_*}{f_*}>\tfrac12\quad(\delta\in I).
    \label{eq:local:spectral-witness}
  \end{equation}
  In particular each $\cL_\delta$ has spectrum inside $\Gamma$ and
  $g_*=\min_{\zeta\in\cG}\re\zeta>0$.
\end{proposition}

The Lean inhabitant is at $k_*=1$, hence $N=2\pi$, with $h=\min(h_0,h_1,\tfrac12)$ and
$I=[1-\eta,1+\eta]$, $\eta=\min(\eta_0,\eta_1,\tfrac12)$, so
$0<\tfrac12\leq\delta_-<1<\delta_+\leq\tfrac32$; the growth floor obeys
\[
  g_*\ \geq\ \tfrac34\bigl(448-4\pi^2\bigr)\ \approx\ 306.4,
  \qquad\text{and in particular}\qquad g_*\geq288 .
\]
These numerals come from the formalisation's choices $\varrho=32$
(Lemma~\ref{lem:local:segment-root}) and $r\leq\tfrac14\re\lambda_*$
(Lemma~\ref{lem:local:isolation}); the paper asserts only $g_*>0$
(\texttt{eq:spectral-gap}).  Three points of shape.  First, the clauses
``$v_h\in\C_c^\infty$ and $\Div_\perp(v_h)_\perp=0$'' are not part of the proposition: the
paper states them with \texttt{eq:smooth-local-velocity} itself, just before the
proposition, and the proposition only says ``for the reference velocity profile on $\R^2$,
$v=v_h$, from equation \eqref{eq:local:smooth-velocity}''.  The Lean record has the same shape: the velocity
is stored as the triple $(h,\Omega,U)$, and the vanishing of $v_h$ outside $D(0,1+h)$ is a
theorem about it (\texttt{parameterizedSmoothedVelocity\_eq\_zero\_of\_radius\_le\_norm}),
not a field; the horizontal divergence-freeness is proved for the assembled global velocity
(Lemma~\ref{lem:global-velocity}), not for $v_h$ on its own.  Second,
the resolvent bound quoted above is the $\Lp^2\to\Hs^2$ one of
Lemma~\ref{lem:local:h2-resolvent}; the paper asks only for the $\Lp^2\to\Hs^1$ bound,
which follows from it, and the deviation is recorded at that node.  Third, two small
differences of the Lean inhabitant: the smoothing width is built as a minimum with
$\tfrac12$, so $h=\tfrac12$ is not excluded where the paper writes
$h\in(0,\tfrac12)$ (hence $h\in(0,\tfrac12]$ above), and $f_*$ is an $\Lp^2$ field with compact support rather than a
$\C_c^\infty$ one (Lemma~\ref{lem:local:witness}); nothing downstream uses either.

\begin{proof}
  \leanok
  \uses{lem:jump-horizontal-mode, lem:local:isolation, lem:local:riesz-perturbation,
    lem:local:witness, lem:local:h2-resolvent, def:local:seed}
  Lemma~\ref{lem:jump-horizontal-mode} supplies $\Omega$, $\kappa$ and $\lambda_*$ at
  $N=2\pi$, and Lemma~\ref{lem:local:isolation} the contour $\Gamma$ with
  $g_*\geq\tfrac34\re\lambda_*$ and the projection $P^0$.
  Lemma~\ref{lem:local:witness} gives $f_*$, $L_*$ with $\re\inner{P^0f_*}{f_*}>\tfrac34$,
  and Lemma~\ref{lem:local:riesz-perturbation} thresholds with
  $\norm{P_\delta-P^0}<\tfrac14$; subtracting gives
  $\re\inner{P_\delta f_*}{f_*}>\tfrac12$, so $P_\delta\neq0$ and the spectrum inside
  $\Gamma$ is non-empty.  Lemma~\ref{lem:local:h2-resolvent} supplies the single constant
  $L$ of \eqref{eq:local:continuation-bound} and identifies the projection of the strong
  generator with the one just estimated.  Assembling the pieces into the record of
  Definition~\ref{def:local:seed} finishes.
\end{proof}
